\begin{helper}[Deletion and reindexing of one Scheide block]
Let a reduced Scheide configuration consist of its base followed by its
ordered two-vertex blocks.  For any selected block index $i$, there is a
tree-sequence $S$ obtained by deleting exactly the two displayed edges
and two displayed vertices of block $i$.  It has length
$\operatorname{len}(T)-2$ and vertex set
\[
 V(S)=V(T)\setminus
 \{u^1_{\gamma_i},u^2_{\gamma_i}\}.
\]
Every base vertex and base edge retains its original index.  Every entry
of every block $r\ne i$ is retained, with its edge immediately before its
new vertex and after the base.  Conversely every edge of $S$ is either a
base edge or an edge of a retained block and is an exact edge
identity of $T$.  No retained edge is incident with either deleted block
vertex.
\end{helper}
\begin{proof}
Put $m=p+2i$, the old index of the first vertex of the selected block.
Define the vertex reindexing map on the new indices by
\[
 r_V(q)=\begin{cases}q,&q<m,\\q+2,&q\ge m,
 \end{cases}
\]
and define the edge reindexing map by
\[
 r_E(q)=\begin{cases}q,&q<m-1,\\q+2,&q\ge m-1.
 \end{cases}
\]
The length equation and position equations in
\noderef{ScheideReducedConfiguration} show that both maps land in the
old index ranges and avoid precisely the two selected positions.  They
are injective: indices on the same side of the gap are translated by the
same amount, while an index below the gap cannot equal an index at least
two positions above it.

Set $S.vertex(q)=T.vertex(r_V(q))$ and
$S.edge(q)=T.edge(r_E(q))$.  Vertex and edge injectivity follow from the
two reindexing maps and the corresponding fields of
\noderef{TreeSequence}.  Since $p\ge2$, the first edge lies before the
deleted block, so the first-edge identity is unchanged.

We verify attachment.  An edge before the gap keeps both its new-vertex
position and its earlier attachment position.  An edge after the gap is
an edge of a block $r>i$.  The block-attachment clause of
\noderef{ScheideReducedConfiguration} says that its endpoints are its
own block vertex and a base vertex.  Both are retained; the block vertex
and its edge are shifted down by two positions, while the base attachment
keeps its original earlier index.  Thus every reindexed edge satisfies
the attachment equation required by \noderef{TreeSequence}.

The structural position equations now give the claimed entry formulas.
All base entries keep their indices.  A block $r<i$ keeps its old
positions, and a block $r>i$ has both positions reduced by two; in either
case its edge remains immediately before its displayed vertex and its
edge index is at least $p-1$.  Euclidean division of an index after the
base by two proves the converse classification: it supplies the unique
retained block and its member $j\in\{0,1\}$.  This also shows directly
that every new edge is one of the old exact edge identities.

The image of $r_V$ consists of every old vertex index except $m,m+1$.
Using the block-position equations and vertex injectivity therefore gives
the asserted vertex-set equality.  Finally an old base edge has endpoints
at base positions, and an old block edge has endpoints at its own block
vertex and a base vertex.  For a retained block its block index differs
from $i$, so vertex injectivity separates that endpoint from both deleted
vertices.  Unfolding \noderef{Incident} proves that no retained edge is
incident with either deleted vertex.
\end{proof}
